\chapter{\ifenglish Background Knowledge and Theory\else ทฤษฎีที่เกี่ยวข้อง\fi}

\section{Cryptography}
Cryptography หรือ วิทยาการรหัสลับ เป็นศาสตร์ที่ว่าด้วยการแปลงข้อความปกติให้กลายเป็นข้อความลับ โดยข้อความลับเป็นข้อความที่ผู้อื่นนอกเหนือจากคู่สนทนาที่ต้องการไม่สามารถเข้าใจได้
โดยทั่วไป กระบวนการปกป้องข้อมูลจะเริ่มต้นก่อนการจัดเก็บหรือการส่งข้อมูล โดยขั้นตอนนี้เรียกว่า การเข้ารหัส (Encryption) 
ซึ่งเป็นการนำข้อมูลต้นฉบับมาผ่านกระบวนการทางคณิตศาสตร์ร่วมกับกุญแจ (Key) ที่เป็นค่าตัวเลขสุ่ม 
เพื่อแปลงเป็นข้อมูลที่ถูกเข้ารหัส (Ciphertext) ในทางกลับกัน เมื่อต้องการอ่านข้อมูลดังกล่าว 
จะต้องนำข้อมูลที่ถูกเข้ารหัสมาผ่านกระบวนการทางคณิตศาสตร์ร่วมกับกุญแจอีกครั้ง เพื่อแปลงกลับเป็นข้อมูลดั้งเดิม 
ซึ่งขั้นตอนนี้เรียกว่า การถอดรหัส (Decryption)

\subsection{Symmetric and Asymmetric Encryption}
การเข้ารหัสในสมัยใหม่ถูกแบ่งออกเป็น 2 ประเภทหลักตามวิธีการใช้กุญแจดังนี้ คือ การเข้ารหัสแบบสมมาตร (Symmetric Encryption) และ การเข้ารหัสแบบอสมมาตร (Asymmetric Encryption)

\subsubsection{Symmetric Encryption}
เป็นกระบวนการเข้ารหัสที่ใช้กุญแจ (Key) ดอกเดียวกันทั้งในขั้นตอนการเข้ารหัสและการถอดรหัส ซึ่งเป็นวิธีที่มีประ-สิทธิภาพสูงและประมวลผลได้อย่างรวดเร็ว ยากที่จะทำการถอดรหัสหากไม่ทราบกุญแจ 
แต่มีข้อจำกัดที่สำคัญคือปัญหาในการกระจายกุญแจ 
เนื่องจากผู้ส่งและผู้รับจำเป็นต้องตกลงและส่งมอบกุญแจผ่านช่องทางที่ปลอด-ภัยล่วงหน้า

\subsubsection{Asymmetric Encryption}
เป็นกระบวนการเข้ารหัสที่ใช้กุญแจสองดอกที่มีความสัมพันธ์กันทางคณิตศาสตร์ที่ยากที่จะแก้ไขแต่ง่ายที่จะตรวจสอบได้
ได้แก่ กุญแจสาธารณะ (Public Key) สำหรับการเข้ารหัส และกุญแจส่วนตัว (Private Key) สำหรับการถอดรหัส 
วิธีการนี้ถูกคิดค้นขึ้นเพื่อแก้ปัญหาการกระจายกุญแจของการเข้ารหัสแบบสมมาตร 
และได้รับการยอมรับอย่างแพร่หลายในปัจจุบัน รวมถึงการประยุกต์ใช้ในระบบการโหวตออนไลน์
แต่ด้วยการที่ต้องใช้กุญแจที่มีขนาดใหญ่มากในการเข้ารหัสเพื่อที่จะทำให้การถอดรหัสทำได้ยาก
ทำให้การประมวลผลช้ากว่าการเข้ารหัสแบบสมมาตร


\subsection{Computational Security and Information-Theoretic Security}

\subsubsection{Computational Security}
Computational Security เป็นความปลอดภัยที่อาศัยสมมติฐานที่ว่าปัญหาทางคณิตศาสตร์บางประเภทมีความยากและซับซ้อนเกินกว่าที่จะแก้ไขได้ภายในระยะเวลาและทรัพยากรที่เหมาะสม 
แนวคิดนี้เป็นรากฐานสำคัญของการเข้ารหัสสมัยใหม่ส่วนใหญ่ ไม่ว่าจะเป็น RSA, ECC, AES 
ตลอดจนอัลกอริทึมการเข้ารหัสยุคหลังควอนตัม (Post-Quantum Cryptography: PQC) ระดับความปลอดภัยของระบบเหล่านี้จะถูกประเมินจากเวลาและทรัพยากรทางคอมพิวเตอร์ที่จำเป็นต้องใช้ในการเจาะระบบ 
อย่างไรก็ตาม เมื่อเทคโนโลยีคอมพิวเตอร์ควอนตัมมีความก้าวหน้ามากขึ้น ระบบการเข้ารหัสเหล่านี้จึงเผชิญกับความเสี่ยงที่จะถูกเจาะทำลาย 
แม้ว่าการเข้ารหัสแบบ PQC จะถูกออกแบบมาเพื่อลดความเสี่ยงดังกล่าว แต่วิธีการเหล่านั้นก็ยังคงพึ่งพาสมมติฐานเชิงคำนวณเป็นหลัก จึงไม่สามารถป้องกันการโจมตีจากระบบที่มีพลังการประมวลผลอย่างไม่จำกัดได้


\subsubsection{Information-Theoretic Security}
Information-Theoretic Security เป็นระดับความปลอดภัยที่ไม่ได้พึ่งพาข้อสมมติฐานด้านขีดความสามารถในการคำนวณใดๆ 
โดยระบบสามารถรับประกันได้ว่าผู้โจมตีจะไม่สามารถถอดรหัสข้อความได้เลย แม้ว่าจะมีพลังการประมวลผลทางคอมพิวเตอร์สูงอย่างไม่จำกัดก็ตาม 
หลักการสำคัญของความปลอดภัยเแบบนี้คือการที่ทำให้รูปแบบการเข้ารหัสยังคงปลอดภัยจากการโจมตีและยังคงรับประกันว่าระดับความปลอดภัยจะไม่เสื่อมถอยลงตามกาลเวลา 
แม้เทคโนโลยีจะพัฒนาไปมากเพียงใดในอนาคต อย่างไรก็ตาม ความท้าทายหลักของระบบนี้คือความซับซ้อนใ่นการนำไปใช้งานจริง
เนื่องจากผู้ใช้งานจำเป็นต้องปฏิบัติตามกฎของโปรโตคอลอย่างเคร่งครัด เพื่อรักษาความสมบูรณ์ของระบบโดยไม่ลดทอนระดับความปลอดภัยลง
ซึ่งการที่จะปลอดภัย\\
นิยามทางคณิตศาสตร์ของความลับที่สมบูรณ์แบบ (Perfect Secrecy) ในระบบนี้ สามารถแสดงได้ด้วยสมการความน่าจะเป็นดังต่อไปนี้:
\[P(M | C) = P(M)\]
โดยที่ $M$ แทนข้อความต้นฉบับ (Plaintext) และ $C$ แทนข้อความที่ถูกเข้ารหัส (Ciphertext) สมการนี้สื่อความหมายว่า ความน่าจะเป็นแบบมีเงื่อนไข (Conditional Probability) ที่ผู้โจมตีจะล่วงรู้ข้อความ $M$ เมื่อดักจับข้อความ $C$ ได้นั้น มีค่าเท่ากับความน่าจะเป็นดั้งเดิมของข้อความ $M$ ซึ่งแปลว่าข้อความที่ถูกเข้ารหัสไม่ได้เปิดเผยข้อมูลหรือเบาะแสใด ๆ เกี่ยวกับข้อความต้นฉบับเพิ่มเติมเลย ~\cite{shannon_secure}

\subsection{One-Time Pad}
การเข้ารหัสแบบวันไทม์แพด (One-Time Pad: OTP)~\cite{vernam_cipher} เป็นขั้นตอนวิธีการเข้ารหัสข้อมูลเพียงรูปแบบเดียวในปัจจุบัน 
ที่ได้รับการพิสูจน์ทางคณิตศาสตร์โดย Claude Shannon ~\cite{shannon_secure} ว่าเป็น Information-Theoretic Security ซึ่งหมายความว่า หากใช้งานอย่างถูกต้อง 
ผู้โจมตีจะไม่สามารถถอดรหัสข้อความไซเฟอร์ (Ciphertext) เพื่อหาข้อความต้นฉบับ (Plaintext) ได้เลย 
ไม่ว่าจะมีพลังประมวลผลหรือเวลามากเพียงใดก็ตามหลักการทำงานพื้นฐานของ OTP ในระบบดิจิทัลมักใช้ตัวดำเนินการทางตรรกศาสตร์ "Exclusive OR" (XOR) 
แบบบิตต่อบิต (Bitwise XOR) ระหว่างบิตของข้อความต้นฉบับ $M$ (Message) และกุญแจ $K$ (Key) 
เพื่อสร้างข้อความไซเฟอร์ $C$ (Ciphertext) ดังสมการ:
\[C = M \oplus K\]
ในฝั่งของผู้รับ เมื่อต้องการถอดรหัส ก็จะนำข้อความไซเฟอร์มาผ่านตัวดำเนินการ XOR ร่วมกับกุญแจชุดเดิมอีกครั้ง 
ก็จะได้ข้อความต้นฉบับกลับคืนมาอย่างสมบูรณ์ ดังสมการ:\[M = C \oplus K\]อย่างไรก็ตาม 
เพื่อให้เป็นไปตามนิยามความปลอดภัยสูงสุด การประยุกต์ใช้งาน OTP ในระบบจริงจำเป็นต้องปฏิบัติตามเงื่อนไขที่เข้มงวด 4 ข้อ
\begin{itemize}
    \item ความยาวของกุญแจต้องเท่ากับหรือยาวกว่าข้อความต้นฉบับ
    \item กุญแจต้องถูกสร้างขึ้นจากตัวเลขสุ่มที่เป็นการสุ่มอย่างแท้จริง
    \item กุญแจต้องถูกรักษาเป็นความลับสุดยอด ระหว่างผู้ส่งและผู้รับเท่านั้น
    \item กุญแจต้องถูกใช้งานเพียงครั้งเดียวแล้วทำลายทิ้ง ห้ามนำกลับมาใช้ซ้ำโดยเด็ดขาด
\end{itemize}

\section{Quantum Key Distribution}
เทคโนโลยีการสื่อสารที่นำหลักการทางกลศาสตร์ควอนตัมมาประยุกต์ใช้ในการสร้างและแลกเปลี่ยน กุญแจเข้ารหัส (Key) ระหว่างผู้ส่งและผู้รับร่วมกันอย่างปลอดภัยสูงสุด
โดยอาศัยหลักการทางกลศาสตร์ควอนตัมเพื่อทำให้การส่งกุญแจเข้ารหัสมีความปลอดภัย
\subsection{Quantum Property}
\subsubsection{Superposition}
Superposition เป็นสมบัติของควอนตัมที่ทำให้สถานะของระบบควอนตัมสามารถอยู่ในสถานะต่างๆ ได้พร้อมกัน
จนกว่าจะมีการวัดค่าเกิดขึ้น คุณสมบัตินี้เป็นกลไกสำคัญที่ทำให้ผู้ส่งสามารถเข้ารหัสข้อมูลลงในอนุภาคโฟตอน (Photon) ได้
\subsubsection{No Cloning Theorem}
หลักการพื้นฐานทางกลศาสตร์ควอนตัม หากไม่รู้ว่าอนุภาคควอนตัมอยู่ในสถานะใด จะไม่สามารถสร้างสำเนาของอนุภาคนั้นได้อย่างสมบูรณ์แบบ
ซึ่งทำให้ผู้ที่ดักฟังข้อมูลไม่สามารถคัดลอกข้อมูลไปได้
\subsubsection{Measurement}
การวัดค่าสถานะทางควอนตัมจะทำให้สถานะของระบบควอนตัมเปลี่ยนแปลงไปสู่สถานะใดสถานะหนึ่ง หรือก็คือการทำให้ Superposition สลายไป
ซึ่งผลลัพธ์ที่ได้จากการวัดค่าสถานะทางควอนตัมจะขึ้นอยู่กับรูปแบบวิธีที่ใช้ในการวัดค่า

\subsection{The BB84 Protocol~\cite{bb84}}
หนึ่งในโปรโตคอลการกระจายกุญแจเข้ารหัสแบบควอนตัม (Quantum Key Distribution)
อาศัยการเข้ารหัสข้อมูลลงบนสถานะโพลาไรเซชัน (Polarization) ของอนุภาคโฟตอนเดี่ยว โดยใช้แกนอ้างอิงหรือ "รูปแบบ (Basis)" ในการวัด 2 รูปแบบที่ตั้งฉากกัน (Conjugate Bases) ได้แก่ รูปแบบแนวตรง ($+$) ซึ่งมีมุม $0^\circ$ และ $90^\circ$ และ รูปแบบแนวทแยง ($\times$) ซึ่งมีมุม $45^\circ$ และ $135^\circ$
กระบวนการทำงานของโปรโตคอล BB84 สามารถแบ่งออกเป็น 5 ขั้นตอนหลัก ดังนี้:
\begin{itemize}
    \item ทำการสุ่มสร้างชุดข้อมูลบิต (0 หรือ 1) และสุ่มเลือกรูปแบบ ($+$ หรือ $\times$) ที่ใช้ในการส่งแล้วส่งโฟตอนเหล่านี้ไปให้ผู้รับผ่านช่องทางสื่อสารควอนตัม
    \item ผู้รับทำการสุ่มเลือกรูปแบบ (Basis) ในการวัดค่าสถานะทางควอนตัมของโฟตอนที่ได้รับของแต่ละบิตซึ่งมีโอกาสถูกเพียง 50\% ดังนั้นผู้รับจะมีข้อมูลที่ถูกต้องเพียง 50\% ของที่ได้รับ
    \item ผู้ส่งและผู้รับทำการนำรูปแบบมาเปรียบเทียบกันเพื่อให้รู้ว่าข้อมูลบิตไหนที่ถูกต้องผ่านช่องทางสื่อสารปกติ
    \item ผู้รับแหละผู้ส่งจะเลือกบิตที่มีการใช้รูปแบบที่ตรงกันเป็นกุญแจเข้ารหัส
    \item ตรวจสอบการดักฟังโดยการสุ่มเลือกบิตมาเปรียบเทียบกันผ่านช่องทางสื่อสารปกติ
\end{itemize}

\section{\ifenglish%
\ifcpe CPE \else ISNE \fi knowledge used, applied, or integrated in this project
\else%
ความรู้ตามหลักสูตรซึ่งถูกนำมาใช้หรือบูรณาการในโครงงาน
\fi
}
\subsubsection{Discrete Mathematics}
นำมาใช้ให้เหตุผลทางคณิตศาสตร์ในการพิสูจน์ความปลอดภัยของโปรโตคอล
\subsubsection{Probability and Statistics}
นำมาใช้ในการวิเคราะห์ความน่าจะเป็นของการดักฟังข้อมูล
\subsubsection{Quantum Computing}
ใช้เป็นความรู้พื้นฐานในการทำความเข้าใจโปรโตคอล และหลักการทำงานของโปรโตคอล
\subsubsection{Theory of Computation and Algorithm}
ใช้เป็นความรู้พื้นฐานในการวิเคราะห์ความซับซ้อนของการเข้ารหัสว่าทำไมถึงยากต่อการคำนวณ และทำความเข้าใจอัลกอริทึมที่ใช้ในการแก้ไขปัญหาของการเข้ารหัส

\section{\ifenglish%
Extracurricular knowledge used, applied, or integrated in this project
\else%
ความรู้นอกหลักสูตรซึ่งถูกนำมาใช้หรือบูรณาการในโครงงาน
\fi
}
\subsubsection{Information Theoretic Security and Computational Security}
ใช้เป็นความรู้พื้นjฐานในการทำความเข้าใจหลักการของการเข้ารหัสว่าทำไมถึงมีความปลอดภัยและจะปลอดภัยได้มากแค่ไหน
นำมาช่วยตัดสินใจเลือกโปรโตคอลที่เหมาะสมกับโครงงานเพื่อความปลอดภัยสูงสุด
\subsubsection{BB84 Protocol}
ใช้เป็นโปรโตคอลในการกระจายกุญแจเข้ารหัสแบบควอนตัมที่ใช้ในโครงงาน
\subsubsection{One-Time Pad}
การเข้ารหัสข้อมูลแบบสมบูรณ์แบบที่ใช้ในการเข้ารหัสข้อมูลในโครงงานเพื่อให้มีความปลอดภัยสูงสุดในระหว่างการส่งข้อมูล